\section{How \cluethreed{} works}
\label{sec:clue3d}

\cluethreed{} is the existing CMS algorithm that groups LayerClusters across
calorimeter layers into Tracksters using the density-based idea of
CLUE~\cite{Rovere2020clue}. The studies use the default HLT \cluethreed{}
configuration without tuning for displaced showers, except where a parameter
variation is explicitly stated.

\subsection{Neighbourhoods projected from the interaction point}

Let LayerCluster $i$ have position $(x_i,y_i,z_i)$, transverse radius
$r_i=\sqrt{x_i^2+y_i^2}$, azimuth $\phi_i$, layer index $\ell_i$, and energy
$E_i$.
To compare a candidate $j$ with $i$, it projects $j$ onto the plane of $i$
along the ray from the origin through $j$. The projected radius is
\begin{equation}
  r_{j\to i}=z_i\frac{r_j}{z_j}.
\end{equation}
The squared transverse distance used by the implementation is then
\begin{equation}
  d_{XY}^2(i,j)=(r_i-r_{j\to i})^2
                    +r_{j\to i}^{\,2}(\Delta\phi_{ij})^2,
  \label{eq:clueprojective}
\end{equation}
where $\Delta\phi_{ij}$ is the azimuthal difference wrapped into
$[-\pi,\pi]$. The second term uses the small-angle arc distance in azimuth.
This is a projected distance, not Euclidean distance of the two LayerClusters. Exchanging $i$ and
$j$ changes the reference plane and can change its value, so $d_{XY}(i,j) \neq d_{XY}(j,i)$.

Candidate searches are restricted to three layers
before and after the reference layer. They also use nearby position-based
$(\eta,\phi)$ buckets called tiles: up to two tiles on either side for density estimation,
and one for the nearest-higher search.

\subsection{Density and Trackster formation}

The local density of LayerCluster $i$ is
\begin{equation}
  \rho_i=E_i+0.2\sum_{j\in\mathcal{N}_i}E_j,
  \label{eq:cluedensity}
\end{equation}
where $\mathcal{N}_i$ contains candidate LayerClusters on other layers with
$d_{XY}(i,j)<\dc$, and $\dc=\SI{1.8}{cm}$. The algorithm then finds the
nearest higher-density LayerCluster $h(i)$ from another layer, minimizing the
projected transverse distance, and records
$\delta_i=d_{XY}(i,h(i))$. This search can find neighbours beyond $\dc$;
if none is found, the distance is treated as infinite.

LayerClusters are classified as follows:
\begin{itemize}
  \item \textbf{Seeds} start new Tracksters. They require
  $\rho_i\geq\rhoc=\SI{0.6}{GeV}$ and separation from nearest higher:
  $\delta_i>\dm=\SI{1.8}{cm}$.
  \item \textbf{Outliers} have low density, $\rho_i<\rhoc$, and a distant
  nearest-higher neighbour, $\delta_i>2\dm=\SI{3.6}{cm}$.
  \item \textbf{Followers} are the remaining LayerClusters and inherit the
  Trackster assignment of their nearest-higher neighbour, if one exists.
\end{itemize}

Starting from each seed, its Trackster index is propagated through the
follower links (Fig.~\ref{fig:rechits-b}). Each assigned LayerCluster belongs
fully to one Trackster.
Outlier status propagates the same way: a follower whose nearest-higher chain
terminates in an outlier instead of seed is itself treated as an outlier.
